CPT conjugation is the composition of grade involution, reversal, and the flag conjugate $v_{k} \leftrightarrow w_{k}$ (Definition \claimref{obj.dark-matter.witt-flag}) on each tower's Clifford bundle; the grade involution negates every vector, the timelike one included (the T; the spatial inversion is the P). Because of the reversal the composite is an anti-automorphism, and it acts on fields as a symmetry does (the convention stated in \claimref{dark-matter.premise-cpt-covariance}): on vectors by its vector part, and on a connection through the Lie algebra automorphism it induces, so that curvature is fixed. In that convention it acts on fields as $C = \mathrm{Ad}(I_{6}) \circ \kappa$, $\kappa$ complex conjugation: its vector part $-1$ is the rotor $I_{6}$, and the flag conjugate that preserves the frozen point's action is complex conjugation, which carries the flag $W$ to $\overline W$, in the other family (\claimref{synthesis.cpt-conjugation-pseudoscalar-rotor-times-complex-conjugation}). So $C$ is antilinear. Structure by structure: the fibre metric is fixed by grade involution, reversal and the exchange of a null pair, which is an isometry of the completion's hyperbolic plane; the pairing on the Fock module sends the grade-$k$ state to the grade-$(3-k)$ state \cite[Prop.~5.6]{ButlerT12} (\claimref{standard-model.chirality-two-halves-spinor-module-conjugation}), which is exactly what the flag conjugate does; and the action is the frozen point's, built from the metric and the completion. The derived gauge sector (Fact \claimref{dark-matter.imported-inputs-standard-model-chain}) is CPT-even by the standard theorem for local Lorentz-invariant actions, but that theorem concerns the antiunitary CPT operator of quantum field theory; that $C$, antilinear and so of Wigner's antiunitary type, realises that operator is argued, not proved, so the gauge sector's CPT-evenness is imported and carried as a premise (\claimref{dark-matter.premise-cpt-covariance}).

So $C$ is an isomorphism of the kinematic data exchanging the flag's two orientations, and a quantity the structures determine would be fixed by it while the orientation is not. The last step reads the bit against Definition \claimref{obj.dark-energy.four-ways-number-enters-theory-shape} directly. Let $S_{\mu}$ be the action functionals, $\mu$ the norm $Q(I) = \Lambda/3$ of the frozen point together with whatever couplings the imported gauge sector carries, and $\mathcal{S}_{\mu}$ their solution spaces. The premise's clause for the action says $S_{\mu} \circ C = S_{\mu}$ for every $\mu$, the gauge couplings included; for the norm this is visible directly, since $C$ sends the real vector $I$ to $-I$, so $Q(CI) = Q(I)$ and $C$ does not move $\mu$, and for the gauge couplings it is the premise. A critical point of a $C$-invariant functional is carried by the involution $C$ to a critical point of the same functional, so $C(\mathcal{S}_{\mu}) = \mathcal{S}_{\mu}$ for every $\mu$. Let $q = \pm 1$ be the orientation bit, a function of the solution (the descent plants its flag once, Definition \claimref{obj.dark-matter.witt-flag}, so $q$ is conserved). By the premise that the flag is a datum (\claimref{dark-matter.premise-flag-as-datum}), the sign of the flag's complex structure $J$ relative to the Fock module's unit $I_7$ (the central pseudoscalar of the odd extension, \claimref{dark-matter.premise-flag-as-datum}) is configuration data of each solution, planted by the descent on the three-pair structure, and not part of the kinematic structure: were it kinematic, the bit would be constant on every solution space of the theory carrying it, hence (D) there, and the two orientations would be two isomorphic theories rather than two solutions of one. Under that premise, $C$ exchanges the two orientations, $q(Cx) = -q(x)$. For every $\mu$ with $\mathcal{S}_{\mu}$ nonempty, which holds for the $\mu$ of our universe, $q$ therefore takes both values on $\mathcal{S}_{\mu}$, so it is not constant on $\mathcal{S}_{\mu}$: $q$ is (I). By Proposition \claimref{dark-energy.four-alternatives-are-exhaustive-three-are} it is then neither (D) nor (P), each of which requires constancy on every $\mathcal{S}_{\mu}$; and it is not (K), since $x$ and $Cx$ have isomorphic kinematic structures, $C$ being that isomorphism, while $q(Cx) \neq q(x)$. The step excludes a parameter of the laws for the bit directly, because a parameter of the laws is constant on each solution space and the bit is not; no exhaustiveness of the kinematic, dynamical and initial-data alternatives enters. The relative-versus-absolute split is the Standard-Model paper's chirality proposition, imported under Hypothesis (S) (\claimref{standard-model.hyp-s-slot-statistics}) (Fact \claimref{dark-matter.imported-derived-generation-flag}); the flag-covariance of the $\nu_{R}$ identification is the exchange of the singlet grades, which the legacy script verifies on the three-pair Fock model. The script of this version tests the logical step and the frozen-point part of the action, not the premises: it applies $C = \mathrm{Ad}(I_{6}) \circ \kappa$ explicitly, in Gaussian rational arithmetic, and confirms that on fixed rational data in the six-dimensional ambient space it is $-1$ on vectors, fixes the curvature bivector, and fixes the norms and the integrand $\epsilon(I, V, F, F)$, that on complexified data it is antilinear and carries each of them to its complex conjugate (with controls in which negating $I$ alone, or a real-linear flag conjugate in place of $\kappa$, flips the integrand), that the flag conjugate, on the split real form, is an isometry of the three null pairs exchanging the two pure vacua $v_{1}v_{2}v_{3}$ and $w_{1}w_{2}w_{3}$, and, exhaustively on small finite theories, that a $C$-odd quantity on $C$-stable solution spaces is (I) and not (K), with a control in which dropping $C$-stability makes it (P).
